% ============================================================
% Capítulo 2 - Referencial Teórico
% Bases: artigo ICEIS (pt-br/main.tex), Montero et al. (2025, ICEIS e Information)
% e Montero et al. (2026). Marcadores [VERIFICAR] indicam dados a confirmar.
% ============================================================
\chapter{Referencial Teórico}
\label{cap:referencial}

Este capítulo apresenta os conceitos em que o trabalho se apoia. A Seção~\ref{sec:ia} situa a inteligência artificial e sua evolução; a Seção~\ref{sec:llm} descreve os LLMs; a Seção~\ref{sec:embedding} apresenta os \textit{embeddings}, a representação vetorial de textos em que a recuperação se apoia; a Seção~\ref{sec:chunking}, a divisão dos documentos em trechos, o \textit{chunking}; e a Seção~\ref{sec:rag}, a arquitetura RAG, que combina esses componentes com um LLM. As duas últimas seções descrevem o contexto de aplicação: a PDSA-RS, com os trabalhos anteriores do grupo, e a PNIP, com a certificação higiênico-sanitária de embarcações.

% ------------------------------------------------------------
\section{Inteligência Artificial}
\label{sec:ia}

A inteligência artificial (IA) é o campo da computação que estuda e constrói sistemas capazes de executar tarefas que, quando feitas por pessoas, exigem inteligência, como raciocinar, reconhecer padrões, compreender a linguagem e tomar decisões \cite{Russell2020}. A pergunta sobre se máquinas podem pensar foi formulada por \citeonline{Turing1950}, que propôs substituí-la por um teste operacional, o jogo da imitação. O termo ``inteligência artificial'' surgiu poucos anos depois, na proposta do seminário de Dartmouth de 1956, que partia da conjectura de que todo aspecto da aprendizagem ou da inteligência poderia ser descrito com precisão suficiente para ser simulado por uma máquina \cite{McCarthy2006}.

\subsection{Da IA simbólica ao aprendizado de máquina}

As primeiras décadas foram dominadas pela IA simbólica, em que o conhecimento é representado explicitamente por regras e símbolos manipulados por algoritmos de busca e inferência lógica. Os sistemas especialistas, que codificavam em regras o conhecimento de especialistas humanos, tiveram aplicação comercial nas décadas de 1970 e 1980, mas mostraram-se caros de construir e manter e frágeis diante de situações não previstas pelas regras \cite{Russell2020}. Essas dificuldades, somadas a expectativas não cumpridas, levaram a períodos de redução de investimento conhecidos como ``invernos da IA''.

O aprendizado de máquina mudou a abordagem: em vez de escrever as regras, o sistema as induz a partir de exemplos \cite{Bishop2006}. No aprendizado supervisionado, o modelo aprende a mapear entradas em saídas a partir de exemplos rotulados; no não supervisionado, procura estrutura em dados sem rótulos; e no aprendizado por reforço, aprende uma política de ações a partir das recompensas que recebe do ambiente \cite{Russell2020}. Em todos os casos, o desempenho passa a depender da quantidade e da qualidade dos dados, e não apenas do conhecimento que o projetista consegue explicitar.

\subsection{Redes neurais e aprendizado profundo}

As redes neurais artificiais são modelos compostos por camadas de unidades que calculam somas ponderadas de suas entradas seguidas de funções de ativação não lineares. O algoritmo de retropropagação do erro \cite{Rumelhart1986} tornou viável treinar redes com camadas intermediárias, ajustando os pesos por gradiente descendente para minimizar uma função de perda. O aprendizado profundo empilha muitas dessas camadas, de modo que cada uma aprende representações progressivamente mais abstratas dos dados, sem que as características precisem ser definidas manualmente \cite{LeCun2015,Goodfellow2016}.

O marco da retomada foi a vitória de uma rede convolucional profunda na competição ImageNet de 2012 \cite{Krizhevsky2012}, possibilitada por grandes conjuntos de dados rotulados e pelo processamento paralelo em GPUs. Para dados sequenciais, como texto, as redes recorrentes, em particular as LSTM \cite{Hochreiter1997}, foram por anos a arquitetura dominante, até a chegada do Transformer \cite{Vaswani2017}, base dos modelos de linguagem descritos a seguir.

% ------------------------------------------------------------
\section{Large Language Models}
\label{sec:llm}

Um modelo de linguagem é um modelo computacional que atribui probabilidades a sequências de palavras ou, de forma equivalente, estima a probabilidade da próxima palavra dado o texto anterior. Os \textit{Large Language Models} (LLMs) são modelos de linguagem baseados no Transformer, com bilhões de parâmetros e pré-treinados em corpora massivos, que adquirem com a escala a capacidade de executar uma grande variedade de tarefas de linguagem natural \cite{Zhao2023}.

\subsection{Dos modelos estatísticos aos neurais}

Os primeiros modelos de linguagem eram estatísticos: os modelos de \textit{n}-gramas estimam a probabilidade de uma palavra a partir da contagem das $n-1$ palavras que a precedem em um corpus. São simples, mas não generalizam para sequências que não aparecem nos dados e não capturam dependências além da janela fixa. \citeonline{Bengio2003} propuseram um modelo de linguagem neural que aprende, junto com a distribuição de probabilidade, uma representação vetorial contínua para cada palavra, de modo que palavras com significados próximos recebem vetores próximos. Essa ideia está na origem dos \textit{embeddings} (Seção~\ref{sec:embedding}). As redes recorrentes e LSTM estenderam o contexto considerado, mas processam o texto um \textit{token} por vez, o que dificulta paralelizar o treinamento e preservar dependências longas \cite{Hochreiter1997}.

\subsection{A arquitetura Transformer}

O Transformer \cite{Vaswani2017} eliminou a recorrência e passou a modelar as relações entre todas as posições de uma sequência por meio do mecanismo de atenção. Na autoatenção, cada \textit{token} é projetado em três vetores, consulta ($Q$), chave ($K$) e valor ($V$), e a saída de cada posição é uma média dos valores de todas as posições, ponderada pela similaridade entre sua consulta e as chaves das demais, como mostra a Equação~\ref{eq:atencao}, em que $d_k$ é a dimensão das chaves.

\begin{equation}
	\label{eq:atencao}
	\mathrm{Atenção}(Q, K, V) = \mathrm{softmax}\!\left(\frac{QK^{\top}}{\sqrt{d_k}}\right) V
\end{equation}

Com isso, o modelo pode relacionar diretamente palavras distantes, como um pronome e o substantivo a que se refere. A atenção com múltiplas cabeças executa várias dessas operações em paralelo, cada uma com projeções próprias, o que permite capturar tipos diferentes de relação. Como a atenção não depende da ordem, informações de posição são somadas às representações de entrada. E como todas as posições são processadas ao mesmo tempo, o treinamento pode ser paralelizado em GPUs, o que tornou viável treinar modelos muito maiores \cite{Vaswani2017}.

O texto não entra no modelo como palavras, mas como \textit{tokens}, unidades obtidas por algoritmos de segmentação em subpalavras, como o \textit{Byte Pair Encoding} \cite{Sennrich2016}. Palavras frequentes tendem a corresponder a um único \textit{token}, e palavras raras são divididas em fragmentos. Por isso, os limites de entrada e de saída dos LLMs, assim como o custo de uso dos modelos hospedados, são medidos em \textit{tokens}, e não em caracteres ou palavras.

\subsection{Pré-treinamento, ajuste e alinhamento}

O Transformer original tinha um codificador e um decodificador. Os modelos que vieram depois usam, em geral, apenas uma das partes. O BERT usa o codificador, pré-treinado para prever palavras mascaradas a partir do contexto à esquerda e à direita, e se tornou a base de modelos de compreensão e de \textit{embeddings} \cite{Devlin2019}. A família GPT usa o decodificador, pré-treinado para prever o próximo \textit{token}, e é ajustada depois para tarefas específicas \cite{Radford2018}. Os LLMs atuais seguem, em sua maioria, esta segunda linha.

Ao escalar esse pré-treinamento, o GPT-3, com 175 bilhões de parâmetros, mostrou que um modelo suficientemente grande executa tarefas novas a partir de apenas uma instrução ou de poucos exemplos inseridos no \textit{prompt}, sem atualizar seus pesos, o que se chamou de aprendizado em contexto \cite{Brown2020}. Modelos pré-treinados em larga escala e adaptáveis a muitas tarefas passaram a ser chamados de modelos de fundação \cite{Bommasani2021}. Um modelo apenas pré-treinado, porém, continua textos, e não necessariamente segue instruções. O ajuste por instruções, combinado com o aprendizado por reforço a partir de feedback humano (RLHF), alinha o comportamento do modelo ao que os usuários pedem: no InstructGPT, as respostas de um modelo de 1,3 bilhão de parâmetros ajustado dessa forma foram preferidas pelos avaliadores às do GPT-3 de 175 bilhões \cite{Ouyang2022}.

\subsection{Modelos abertos, execução local e multimodalidade}

Os primeiros LLMs de alto desempenho eram proprietários e acessíveis apenas por APIs. A família LLaMA mostrou que modelos com pesos abertos, treinados apenas com dados públicos, podiam competir com modelos fechados muito maiores \cite{Touvron2023}. Modelos abertos podem ser executados na infraestrutura da própria instituição, de modo que os documentos processados não saem dela, o que é relevante para órgãos públicos que tratam dados pessoais e documentos de terceiros \cite{Brasil2018}. Para caber em \textit{hardware} modesto, esses modelos costumam ser quantizados, isto é, ter seus pesos representados com menos bits, por exemplo 4 em vez de 16, com perda pequena de qualidade \cite{Dettmers2023}. Ferramentas como o Ollama, usado nos trabalhos do grupo e neste, simplificam a execução local desses modelos.

Os LLMs também deixaram de ser apenas textuais. Os modelos de visão e linguagem combinam um codificador de imagens com um LLM, de modo que imagens e texto possam ser fornecidos na mesma entrada, e são ajustados com instruções que envolvem imagens \cite{Liu2023}. Esses modelos permitem, por exemplo, pedir que se descreva uma fotografia ou que se verifique se ela mostra determinado equipamento, o que interessa a processos cujas evidências incluem registros fotográficos, como a certificação HS.

\subsection{Parâmetros de geração}

A cada passo, o LLM produz uma distribuição de probabilidade sobre o vocabulário, e a estratégia de decodificação define qual \textit{token} é escolhido. A decodificação gulosa escolhe sempre o mais provável, e a amostragem sorteia o próximo \textit{token} da distribuição. A temperatura reescala essa distribuição antes do sorteio: valores baixos a concentram nos \textit{tokens} mais prováveis e tornam a saída mais determinística, e valores altos a achatam e aumentam a diversidade. A amostragem por núcleo restringe o sorteio ao menor conjunto de \textit{tokens} cuja probabilidade acumulada atinge um limiar, evitando tanto a repetição em laço, comum na decodificação gulosa, quanto palavras improváveis \cite{Holtzman2020}. Em tarefas administrativas, em que se espera consistência, costuma-se usar temperaturas baixas \cite{Montero2026}.

Outro parâmetro é a janela de contexto, o número máximo de \textit{tokens} que o modelo considera de uma vez, somando a entrada e a saída. Janelas maiores permitem incluir documentos inteiros no \textit{prompt}, mas não garantem que o modelo use bem todo o conteúdo: o desempenho tende a cair quando a informação relevante está no meio de um contexto longo, e não em seu início ou fim \cite{Liu2024}.

\subsection{Limitações em domínios especializados}

Apesar da capacidade de generalização, os LLMs têm limitações que pesam em domínios regulados. A primeira é o corte de conhecimento: o modelo só conhece o que estava em seus dados de treinamento, de modo que normas publicadas ou alteradas depois disso lhe são desconhecidas. A segunda é a cobertura: normas locais e técnicas, como portarias de um ministério brasileiro, são raras nos corpora de pré-treinamento, e o desempenho dos modelos cai em domínios distantes desses dados \cite{Gururangan2020}. A terceira é a alucinação, a produção de conteúdo fluente e plausível, mas sem apoio nos dados ou contrário a eles \cite{Ji2023}. Em um parecer técnico, isso equivale a citar um requisito inexistente ou uma evidência que o processo não contém. Por fim, as respostas não indicam sua origem, o que dificulta verificá-las.

Duas estratégias tratam dessas limitações: especializar o próprio modelo, por \textit{fine-tuning}, isto é, continuando seu treinamento em dados do domínio \cite{Gururangan2020}, ou fornecer o conhecimento no momento da consulta, por recuperação. Para incorporar conhecimento factual, a recuperação superou o \textit{fine-tuning} não supervisionado \cite{Ovadia2024}, e é a estratégia seguida neste trabalho, como nos trabalhos anteriores do grupo \cite{Montero2025,Montero2026}: a norma da certificação HS mudou duas vezes desde 2020, os pareceres precisam indicar o requisito em que se baseiam, e não há pareceres rotulados em número suficiente para treinar um modelo. Antes de tratar da recuperação, a próxima seção apresenta os \textit{embeddings}, a representação de textos em que ela se apoia.

% ------------------------------------------------------------
\section{Embeddings}
\label{sec:embedding}

Um \textit{embedding} é uma representação de um texto (uma palavra, uma frase ou um trecho inteiro) como um vetor de números reais de dimensão fixa, construída de modo que textos com significados próximos fiquem próximos no espaço vetorial. É essa propriedade que permite encontrar trechos relevantes para uma consulta mesmo quando eles não compartilham palavras com ela, e é sobre ela que se constrói a recuperação dos sistemas RAG (Seção~\ref{sec:rag}).

\subsection{Representações esparsas e densas}

As representações clássicas da recuperação de informação são esparsas: cada documento é um vetor com uma dimensão por termo do vocabulário, quase todas nulas, ponderadas por esquemas como o TF-IDF, que valoriza termos frequentes no documento e raros na coleção \cite{Manning2008}. O BM25, um refinamento probabilístico dessa ideia, continua sendo uma linha de base forte \cite{Robertson2009}. Essas representações, porém, só reconhecem correspondências exatas de termos: uma consulta sobre ``potabilidade da água'' não encontra um trecho que fale em ``água própria para consumo humano''.

As representações densas resolvem esse descompasso de vocabulário. O word2vec aprende, a partir das palavras que costumam aparecer juntas, vetores de algumas centenas de dimensões em que relações semânticas viram relações geométricas; o exemplo clássico é que o vetor de ``rei'' menos o de ``homem'' mais o de ``mulher'' fica próximo do de ``rainha'' \cite{Mikolov2013}. Esses vetores, porém, são estáticos: uma palavra tem o mesmo vetor em qualquer contexto. Os modelos baseados no Transformer, como o BERT, produzem representações contextuais, em que o vetor de cada \textit{token} depende da frase em que ele aparece \cite{Devlin2019}.

\subsection{Embeddings de sentenças e de trechos}

Para recuperar trechos, é preciso representar um texto inteiro em um único vetor. O Sentence-BERT ajusta o BERT em uma arquitetura siamesa, em que dois textos passam pelo mesmo codificador e o treinamento aproxima os vetores de pares semanticamente equivalentes, o que permite comparar textos pela similaridade de cosseno de seus vetores \cite{Reimers2019}. O \textit{Dense Passage Retrieval} (DPR) usa dois codificadores BERT, um para as perguntas e outro para os trechos, treinados para que cada pergunta fique próxima dos trechos que a respondem; em perguntas de domínio aberto, superou o BM25 em 9 a 19 pontos percentuais de acurácia na recuperação dos 20 primeiros trechos \cite{Karpukhin2020}. É esse recuperador que compõe o RAG original, descrito na Seção~\ref{sec:rag} \cite{Lewis2020}.

Como os vetores dos documentos são calculados uma vez, na indexação, e a consulta é codificada separadamente, esses modelos são chamados de bicodificadores. A alternativa, os codificadores cruzados, processa a consulta e o trecho juntos e estima melhor a relevância, mas precisa ser executada para cada par, o que só é viável para reordenar um pequeno conjunto já recuperado; é esse o papel do reranqueamento nos sistemas RAG avançados (Seção~\ref{sec:rag}) \cite{Gao2023}.

A escolha do modelo de \textit{embedding} importa em domínios como o deste trabalho, em que os documentos estão em português. O BGE-M3 é um modelo multilíngue que aceita entradas de até 8.192 \textit{tokens} e suporta, com o mesmo modelo, recuperação densa, esparsa e por múltiplos vetores \cite{Chen2024}. Trabalhos anteriores do grupo usaram o modelo \texttt{bge-large-en-v1.5} \cite{Montero2025}; neste trabalho foram usados o BGE-M3, executado localmente, e o \texttt{gemini-embedding-001}, do Google, acessado por API (Capítulo~\ref{cap:estudo-caso}).

\subsection{Similaridade e busca vetorial}

A proximidade entre dois vetores é medida, em geral, pela similaridade de cosseno, que considera apenas o ângulo entre eles, e não seu comprimento (Equação~\ref{eq:cosseno}) \cite{Manning2008}. Ela vai de $-1$ a $1$, e vetores que apontam na mesma direção têm similaridade 1. Quando os vetores são normalizados, ela coincide com o produto interno.

\begin{equation}
	\label{eq:cosseno}
	\mathrm{sim}(\mathbf{q}, \mathbf{d}) = \cos\theta = \frac{\mathbf{q} \cdot \mathbf{d}}{\lVert \mathbf{q} \rVert \, \lVert \mathbf{d} \rVert} = \frac{\sum_{i=1}^{n} q_i d_i}{\sqrt{\sum_{i=1}^{n} q_i^2} \, \sqrt{\sum_{i=1}^{n} d_i^2}}
\end{equation}

Encontrar os $k$ vetores mais similares comparando a consulta com todos os vetores armazenados tem custo proporcional ao tamanho da coleção. Os bancos vetoriais usam, por isso, índices de busca aproximada de vizinhos mais próximos, que trocam uma pequena perda de exatidão por buscas muito mais rápidas. O HNSW, um dos mais usados, organiza os vetores em um grafo de várias camadas, em que as superiores têm poucas conexões longas e as inferiores muitas conexões curtas; a busca começa no topo e desce refinando a vizinhança \cite{Malkov2020}. O pgvector acrescenta ao PostgreSQL um tipo de dado vetorial, operadores de distância e índices HNSW, de modo que os vetores fiquem no mesmo banco relacional que os metadados e possam ser filtrados por eles na mesma consulta \cite{pgvector}. Essa combinação foi usada na PDSA-RS \cite{Montero2025} e é a usada pelo DocFlow.

% ------------------------------------------------------------
\section{Chunking}
\label{sec:chunking}

\textit{Chunking} é a divisão dos documentos em trechos menores antes da indexação. Ele é necessário porque os modelos de \textit{embedding} têm limite de entrada, porque o vetor de um texto longo mistura muitos assuntos e representa mal cada um deles, e porque a janela de contexto do LLM e o custo da inferência limitam quanto texto pode ser inserido no \textit{prompt}. A forma de dividir determina a unidade que será recuperada e citada e, por isso, afeta diretamente a qualidade da recuperação nos sistemas RAG (Seção~\ref{sec:rag}) \cite{Gao2023}.

O tamanho dos trechos envolve uma troca. Trechos pequenos geram vetores mais focados e uma recuperação mais precisa, mas podem perder o contexto necessário para interpretá-los, como a condição enunciada no parágrafo anterior. Trechos grandes preservam o contexto, mas diluem o vetor e levam ao \textit{prompt} conteúdo irrelevante. A sobreposição, em que cada trecho repete o final do anterior, reduz o risco de que uma informação fique partida entre dois trechos, à custa de redundância no índice.

\subsection{Estratégias}

As estratégias mais comuns são:

\begin{itemize}
	\item \textbf{Tamanho fixo}: o texto é cortado a cada $n$ caracteres ou \textit{tokens}, com ou sem sobreposição. É simples e previsível, mas ignora a estrutura do texto e pode cortar frases e artigos ao meio. Na PDSA-RS, por exemplo, foram usados trechos de 1.024 caracteres com sobreposição de 200 \cite{Montero2025}.
	\item \textbf{Recursiva}: o texto é dividido pelo separador de maior nível, como a quebra de parágrafo, e os trechos que continuam acima do limite são divididos de novo pelo separador seguinte, como a quebra de linha, a frase e a palavra. Assim os cortes tendem a cair nos limites naturais do texto, respeitando um tamanho máximo. É a estratégia padrão de bibliotecas como LangChain e LlamaIndex e do DocFlow.
	\item \textbf{Semântica}: o texto é dividido onde a similaridade entre os \textit{embeddings} de frases consecutivas cai, o que indica mudança de assunto.
	\item \textbf{Estrutural}: o texto é dividido segundo a estrutura do próprio documento, como seções, títulos, tabelas ou, em textos legais, artigos.
\end{itemize}

\subsection{Chunking de textos normativos}

Os textos normativos têm uma estrutura explícita que favorece a divisão estrutural. Na legislação brasileira, a Lei Complementar nº~95/1998 estabelece que a unidade básica de articulação dos textos legais é o artigo, desdobrado em parágrafos, incisos, alíneas e itens \cite{Brasil1998}. Cada artigo tende, portanto, a tratar de um único assunto e é a unidade pela qual as normas são citadas. Um trecho que corresponde a um artigo pode ser recuperado e citado como ``Art.~8 da Portaria 310/2020'', ao passo que uma janela de tamanho fixo pode começar no meio de um artigo e terminar no seguinte, perdendo a referência de sua origem.

A literatura confirma que essa escolha importa. Em documentos governamentais tailandeses, \citeonline{Boonmee2026} compararam a divisão recursiva de tamanho fixo, em três tamanhos, com a divisão semântica usando o BGE-M3, e mostraram que o melhor método depende da estrutura de cada tipo de documento. Em normas indonésias, um \textit{framework} que respeita os limites dos artigos (\textit{pasal}) e expande a consulta atingiu 90\% de acurácia nas respostas \cite{Fiddina2026}, e a arquitetura de \citeonline{Montero2026} também segmenta a legislação por artigo. Neste trabalho, o corpus normativo foi segmentado por artigo, com a divisão recursiva apenas como alternativa para artigos longos e textos sem artigos, e cada trecho foi armazenado com o arquivo e o artigo de origem como metadados (Capítulo~\ref{cap:metodologia}).

% ------------------------------------------------------------
\section{Retrieval-Augmented Generation}
\label{sec:rag}

A \textit{Retrieval-Augmented Generation} (RAG) combina um LLM com um mecanismo de recuperação de informação: antes de gerar a resposta, o sistema busca em uma base de conhecimento externa os trechos mais relevantes para a consulta e os insere no \textit{prompt} \cite{Lewis2020,Gao2023}. \citeonline{Lewis2020} descreveram essa combinação como a união de uma memória paramétrica, o conhecimento guardado nos pesos do modelo gerador, com uma memória não paramétrica, um índice vetorial de documentos acessado pelo recuperador denso DPR (Seção~\ref{sec:embedding}) \cite{Karpukhin2020}. Em tarefas de perguntas e respostas que exigem conhecimento factual, o modelo com recuperação superou modelos apenas paramétricos e gerou respostas mais específicas e factuais \cite{Lewis2020}.

\subsection{Funcionamento}

Um sistema RAG opera em duas fases, ilustradas na Figura~\ref{fig:rag}. Na fase de \textbf{indexação}, executada antes das consultas, os documentos da base de conhecimento têm seu texto extraído, são divididos em trechos (\textit{chunks}, Seção~\ref{sec:chunking}), cada trecho é convertido em um vetor por um modelo de \textit{embedding} (Seção~\ref{sec:embedding}), e os vetores são armazenados em um banco de dados vetorial junto com metadados, como o documento de origem. Na fase de \textbf{consulta}, a pergunta do usuário é convertida em vetor pelo mesmo modelo, o banco retorna os $k$ trechos mais similares, e esses trechos são inseridos no \textit{prompt} junto com a pergunta e as instruções. O LLM gera então a resposta condicionada a esse contexto.

\begin{figure}[htb]
	\centering
	\caption{Fases de indexação e de consulta de um sistema RAG.}
	\label{fig:rag}
	\resizebox{\textwidth}{!}{%
	\begin{tikzpicture}[
		font=\small,
		caixa/.style={draw, rounded corners=2pt, align=center, minimum height=10mm, text width=24mm, fill=white},
		banco/.style={draw, cylinder, shape border rotate=90, aspect=0.25, align=center, minimum height=14mm, text width=20mm, fill=black!8},
		>=stealth, every path/.style={->, thick}]
		% indexação
		\node[font=\small\bfseries, anchor=west] at (-1.3,1.4) {Indexação};
		\node[caixa] (docs) at (0,0) {Documentos\\(normas)};
		\node[caixa, right=6mm of docs] (chunk) {Segmentação\\em \textit{chunks}};
		\node[caixa, right=6mm of chunk] (emb1) {Modelo de\\\textit{embedding}};
		\node[banco, right=8mm of emb1] (vdb) {Banco\\vetorial};
		\draw (docs) -- (chunk);
		\draw (chunk) -- (emb1);
		\draw (emb1) -- (vdb);
		% consulta
		\node[font=\small\bfseries, anchor=west] at (-1.3,-1.9) {Consulta};
		\node[caixa] (q) at (0,-3.3) {Consulta\\do usuário};
		\node[caixa, right=6mm of q] (emb2) {Modelo de\\\textit{embedding}};
		\node[caixa, right=6mm of emb2] (busca) {Busca dos $k$\\mais similares};
		\node[caixa, right=8mm of busca, text width=32mm] (prompt) {\textit{Prompt}: instruções,\\consulta e trechos};
		\node[caixa, right=6mm of prompt] (llm) {LLM};
		\node[caixa, right=6mm of llm, text width=18mm] (resp) {Resposta\\com citações};
		\draw (q) -- (emb2);
		\draw (emb2) -- (busca);
		\draw (vdb.south) |- ++(0,-6mm) -| (busca.north);
		\draw (busca) -- (prompt);
		\draw (prompt) -- (llm);
		\draw (llm) -- (resp);
		\draw[dashed] (q.south) -- ++(0,-5mm) -| (prompt.south);
	\end{tikzpicture}}
	\fonte{Elaborado pelo autor, com base em \citeonline{Lewis2020} e \citeonline{Gao2023}.}
\end{figure}

Como o conhecimento fica fora do modelo, atualizar o sistema é reindexar documentos, sem retreinar nada: quando uma portaria é alterada, basta indexar o texto consolidado. E como cada trecho recuperado carrega sua origem, pode-se pedir ao modelo que cite a fonte de cada afirmação, o que permite ao usuário conferir a resposta no documento original. Essas duas propriedades, atualização e rastreabilidade, tornam o RAG adequado a sistemas governamentais \cite{Montero2025,Montero2026}.

\subsection{Paradigmas e variações}

\citeonline{Gao2023} organizam a evolução dos sistemas RAG em três paradigmas. O \textbf{RAG ingênuo} segue apenas o fluxo de indexar, recuperar e gerar descrito acima. O \textbf{RAG avançado} acrescenta etapas antes e depois da recuperação. Antes, a consulta pode ser reescrita, expandida com sinônimos e siglas ou decomposta em subconsultas; depois, os trechos recuperados podem ser reordenados por um modelo de reranqueamento, em geral um codificador cruzado (Seção~\ref{sec:embedding}), ou comprimidos, para que só o conteúdo mais relevante chegue ao LLM. O \textbf{RAG modular} trata o sistema como um conjunto de módulos combináveis, como buscas híbridas, recuperação iterativa e roteamento entre bases, e inclui as arquiteturas em que agentes decidem quando e o que recuperar.

Outra decisão de projeto é a origem da consulta. Em assistentes de perguntas e respostas, a consulta é a própria pergunta do usuário. Em tarefas de análise de processos, como as tratadas neste trabalho, não há uma pergunta: a consulta precisa ser construída a partir da tarefa ou do conteúdo do processo. Uma consulta fixa, igual para todos os processos, recupera sempre os mesmos trechos, enquanto consultas derivadas dos documentos de cada processo fazem o contexto acompanhar o caso. O Capítulo~\ref{cap:metodologia} descreve como o DocFlow constrói essas consultas.

\subsection{Limitações e avaliação}

O RAG transfere parte do problema para a recuperação. Se os trechos relevantes não forem recuperados, o modelo responde sem eles ou com trechos irrelevantes, que podem desviá-lo. Na PDSA-RS, a maior parte das falhas do sistema RAG veio de erros de recuperação \cite{Montero2025}. Mesmo quando os trechos certos são recuperados, nada garante que o modelo os use corretamente: o contexto pode ser longo, com a informação relevante no meio \cite{Liu2024}, ou a norma pode não conter o critério de que a decisão depende.

Por isso, a avaliação de sistemas RAG costuma separar os dois componentes. A recuperação é avaliada com métricas da área de recuperação de informação, como a revocação nos $k$ primeiros resultados (Recall@$k$), que mede se os trechos relevantes estão entre os recuperados \cite{Manning2008}. A geração é avaliada pela correção da resposta em relação a um gabarito ou por critérios como a fidelidade da resposta ao contexto recuperado e sua relevância para a pergunta, que podem ser estimados automaticamente com o auxílio de um LLM, como no \textit{framework} RAGAS \cite{Es2024}. Em tarefas de decisão, como a deste trabalho, a referência mais direta é a decisão tomada pelos especialistas humanos no mesmo processo.

% ------------------------------------------------------------
\section{PDSA-RS e Trabalhos Anteriores}
\label{sec:pdsa}

Este trabalho dá continuidade a uma linha de pesquisa do LUMAC sobre o uso de LLMs e RAG em plataformas governamentais de defesa sanitária, iniciada na Plataforma de Defesa Sanitária Animal do Rio Grande do Sul (PDSA-RS). Criada em 2019, a PDSA-RS apoia a sanidade e a produção animal no estado por meio de uma plataforma digital para a gestão de certificações e o acompanhamento da conformidade sanitária \cite{Montero2025}. Foi desenvolvida pela UFSM em parceria com o Ministério da Agricultura e Pecuária (MAPA), com apoio do Fundo de Desenvolvimento e Defesa Sanitária Animal (FUNDESA), e integra as etapas da certificação de estabelecimentos avícolas e suínos, ligando inspeções veterinárias, laboratórios e autoridades de defesa agropecuária. Produtores e fiscais acompanham pela plataforma os exames sanitários, o processamento das amostras e a emissão dos certificados exigidos para a movimentação de animais no país e no exterior \cite{Montero2025}. A plataforma adota uma arquitetura orientada a microsserviços, com portais para cada público, como o Serviço Veterinário Estadual, os responsáveis técnicos, os laboratórios e o MAPA \cite{Montero2025}.

Sobre essa plataforma, o grupo aplicou RAG pela primeira vez no módulo de certificação avícola: as instruções normativas que definem os exames exigidos foram indexadas, e um LLM pequeno, executado localmente, passou a responder sobre a condição sanitária das solicitações, com acurácia de 55\% contra 12\% do mesmo modelo sem recuperação \cite{Montero2025iceis,Montero2025}. Em seguida, \citeonline{Montero2026} propuseram uma arquitetura genérica de assistente para analistas do setor público e avaliaram seu motor de geração em pareceres sobre recursos administrativos julgados pela JAP da SEAPI-RS. Essa arquitetura foi depois implementada como a plataforma DocFlow, descrita no Capítulo~\ref{cap:metodologia}. Esses trabalhos são analisados em detalhe na Seção~\ref{sec:tr-grupo}.

% ------------------------------------------------------------
\section{PNIP e a Certificação Higiênico-Sanitária}
\label{sec:pnip}

\subsection{A Plataforma Nacional da Indústria do Pescado}

A Plataforma Nacional da Indústria do Pescado (PNIP), instituída pela Portaria MPA nº~361/2024 \cite{MPA2024}, é a plataforma digital oficial do Ministério da Pesca e Aquicultura para o envio, a análise e a emissão das certificações do setor pesqueiro.
Ao centralizar esses fluxos, a plataforma apoia a rastreabilidade que liga cada produto pesqueiro à sua origem legal \cite{Olsen2018}, um dos pilares do combate à pesca ilegal, não declarada e não regulamentada, que responde por uma parcela relevante das capturas mundiais \cite{Agnew2009}.

Cada solicitação é analisada por perfis hierárquicos (analista técnico, coordenador e diretor), o que exige conhecimento aprofundado de normas que mudam com o tempo. A plataforma também se integra a sistemas externos do governo e de parceiros, como SERPRO, CRMV, SISRGP, PesqBrasil e SEI, para validar dados e tramitar documentos.

\subsection{A certificação higiênico-sanitária de embarcações}

A certificação higiênico-sanitária de embarcações pesqueiras é regida pela Portaria SAP/MAPA nº~310/2020, alterada pelas Portarias MPA nº~171/2023 e nº~742/2026 \cite{Portaria310,Portaria171,Portaria742}, e resulta no Certificado Oficial de Boas Práticas Higiênico-Sanitárias a Bordo, válido por cinco anos. A norma estabelece os requisitos que a embarcação deve cumprir, como o programa de autocontrole, a potabilidade da água usada a bordo, o treinamento da tripulação, o controle de pragas e as características do material de revestimento das áreas em contato com o pescado, e traz no Anexo~I a lista de verificação usada na análise. A Portaria 742/2026 acrescentou requisitos específicos para a secagem de bexiga natatória e o congelamento de camarão e lagostim a bordo.

Para solicitar o certificado, o técnico responsável pela embarcação preenche na PNIP a lista de verificação do Anexo~I e anexa as evidências de cada requisito, organizadas por categoria:

\begin{itemize}
	\item programa de autocontrole (PAC) e registros de controle a bordo;
	\item registros de treinamento da tripulação;
	\item documentos do controle de pragas;
	\item laudo de potabilidade da água e declaração do fornecedor de água;
	\item documento do material de revestimento;
	\item declaração de responsabilidade;
	\item registros fotográficos datados da embarcação e de seus equipamentos;
	\item outros documentos.
\end{itemize}

\subsection{A análise técnica e o parecer}

O analista verifica cada item da lista de verificação contra os anexos, registra as não conformidades e redige um parecer técnico, que avalia os dados da embarcação, cada categoria de documentos e cada item da lista e conclui de forma favorável ou desfavorável. Se faltar algum item ou a evidência for insuficiente, o analista solicita complementos, e o técnico tem 30 dias corridos para apresentá-los; caso contrário, a solicitação é indeferida. A decisão depende, portanto, de o conteúdo de cada anexo comprovar o requisito correspondente, e não apenas de o anexo existir. Nos registros fotográficos, por exemplo, o analista verifica se o conjunto de fotos cobre as vistas exigidas da embarcação e se as fotos pertencem a ela.

É esse julgamento que este trabalho procura apoiar. Em comparação com a certificação avícola da PDSA-RS, em que a decisão decorre de dados estruturados de exames laboratoriais, a certificação HS envolve documentos heterogêneos, muitas vezes digitalizados, e dezenas de fotografias por processo, e cabe ao analista julgar se a evidência é suficiente. A revisão apresentada no Capítulo~\ref{cap:trabalhos-relacionados} não encontrou estudos que combinassem os documentos de um dossiê, seus registros fotográficos e as normas recuperadas para apoiar esse tipo de julgamento.
